\section*{RESPONSABILIDADES}

\subsection*{Responsabilidad del Jefe de Grupo}
\begin{enumerate}
    \item Declarar el inicio de una emergencia y su clasificación.
    \item Establecer las funciones de cada uno de los participantes en las operaciones de emergencia.
    \item Coordinar las acciones del grupo y tomar las decisiones que sean necesarias para mitigar las consecuencias del accidente.
    \item Dirigir personalmente las labores de búsqueda o recuperación de Fuente.
\end{enumerate}

\subsection*{Responsabilidades del Responsable de Comunicación}
\begin{enumerate}
    \item Servir de enlace entre el personal que atiende la emergencia y el personal de apoyo localizado en las Oficinas Generales, CNSNS o autoridades.
    \item Recabar toda la información que sea necesaria para:
    \begin{enumerate}[label=\roman*.]
        \item Evaluar la emergencia.
        \item Efectuar cálculos dosimétricos.
        \item Reconstruir el escenario.
    \end{enumerate}
    \item Podrá también cooperar en las labores de búsqueda y recuperación si así lo solicita el Jefe de Grupo aunque sin olvidar su responsabilidad principal.
\end{enumerate}

\subsection*{Responsabilidades de las Brigadas}
\begin{enumerate}
    \item Seguir las indicaciones del Jefe de Grupo.
    \item Cumplir en todo momento las disposiciones del Manual de Seguridad Radiológica.
    \item Reportarse periódicamente para informar e informarse de los avances de la emergencia ante el Jefe de Grupo.
    \item Llevar un estricto control y vigilancia sobre las dosis a fin de no rebasar los límites autorizados.
\end{enumerate}
